\documentclass[a4paper,12pt]{article}
\usepackage{amsmath}
\usepackage[utf8]{inputenc}
\usepackage[pdflatex]{graphicx}
\usepackage{float}

\begin{document}
\title{Artículo Dummy \\ \large{Subtítulo por debajo}}
\author{Gutierrez-Ariza, Carlos M.}
\maketitle

\begin{abstract}
  https://github.com/Charlie-Flin/Proyecto\_Final \\
  Y lo demás es latín de relleno hasta el final
\end{abstract}

\vspace{1cm}
{\bf Keywords::} latín, lore, ipsum, github

\pagebreak

\section{Estado del Arte}
Lorem ipsum dolor sit amet consectetur adipiscing elit montes, maecenas fringilla augue nam penatibus quam egestas, parturient sociis id netus posuere dui magnis. Ac vulputate eros gravida malesuada blandit auctor, senectus magna augue fermentum venenatis elementum, tempus ligula vehicula tristique aptent. Varius justo ornare urna viverra in taciti montes vestibulum, vehicula est ut leo mi dictum orci, torquent gravida tincidunt eget vitae suspendisse lacinia. Cursus at nisl pellentesque luctus faucibus habitasse orci vehicula molestie suscipit senectus, ad primis cum dictum imperdiet magna quisque id nullam.

Gravida malesuada aptent nibh lacinia est mauris, primis luctus imperdiet nullam ultrices, curae risus tortor dictumst integer. A nullam duis conubia tincidunt sodales consequat elementum aenean, vestibulum suspendisse nisi lectus morbi metus sem torquent, ligula magnis vulputate dignissim sollicitudin sociosqu molestie.

Dui elementum ligula interdum sapien taciti maecenas ridiculus aenean arcu, integer porta diam risus phasellus mattis pulvinar mollis malesuada sagittis, penatibus netus rhoncus ullamcorper ut cras at nunc. Et fames sagittis porta lacinia eget eu natoque curabitur rhoncus magna fermentum vestibulum dignissim, ullamcorper fusce nisi orci fringilla ultricies nullam nascetur vel ultrices bibendum euismod. Imperdiet urna lacinia morbi ac placerat mi massa sodales pretium, felis sed parturient ad lobortis inceptos purus tortor laoreet, nec libero netus nunc sagittis nibh sapien nisi.

Urna hac elementum at erat dictumst placerat velit penatibus inceptos, senectus quam purus praesent fringilla rhoncus posuere nullam eu, sodales proin pharetra odio mus nam conubia risus. Eleifend dapibus justo senectus convallis aliquet fringilla ultrices fusce quis facilisis cum, imperdiet sagittis sem magnis curae sociis dictumst mi vestibulum integer. Inceptos laoreet ridiculus varius luctus suspendisse mus senectus ut neque aliquam, nisi eleifend eu sapien placerat turpis ultrices mi tempus nec quis, ligula faucibus lacus lectus ullamcorper montes diam arcu sociis. Ultrices tortor purus luctus ligula blandit gravida non ut, sapien magna nunc fusce a ultricies inceptos, dis sociosqu placerat bibendum fringilla orci metus.

\section{Imagen y Tabla}
\begin{figure}[H]
\includegraphics{prueba.png}
\caption{Pie de figura}
\end{figure}

\begin{center}
    \begin{tabular}{ | l | l | l | l |}
    \hline
    1 & 2 & 3 & 4 \\ \hline
    5 & 6 & 7 & 8 \\ \hline
    9 & 10 & 11 & 12 \\ \hline
    13 & 14 & 15 & 16 \\\hline
    \end{tabular}
\end{center}

\section{aqui una formula}
Formula de ejemplo:
\begin{equation}
\frac{\partial\left(\rho s\right)}{\partial t}=-div\left(J_{s}\right)+\sigma\,\,\,\,\,\,\text{con}\,\,\,\,\,\,\sigma>0\label{eq:3}
\end{equation}
y con esto y bibliografia terminamos

\pagebreak

\nocite{*}
\bibliography{bibliografia}
\bibliographystyle{plain}

\end{document}